\chapter{Evaluation und Testing von LLM-Applikationen}
\label{chap:evaluation}

% ============================================================
% AUTOR-NOTIZ STATT LERNZIELE
% ============================================================
\autornotiz{
2024: SupportPilot ging ohne Evaluation live. Prompt-v1 "fühlte sich gut an" --- aber wir wussten nicht, ob er besser oder schlechter war als ein Münzwurf.
Erst Golden Set (80 Cases, stratifiziert aus echten User-Queries) + Multi-Judge (GPT-4o, Claude 3.5 Sonnet, Gemini 1.5 Pro) + CI-Gate (Pass Rate $\geq$ 85\%, Faithfulness $\geq$ 0.8, Cost $\leq$ €0.02) machte das System steuerbar.
Ergebnis nach 3 Monaten:
- Pass Rate: 72\% $\to$ 94\% (Golden Set, n=80, Agreement $\geq$ 66\%)
- Halluzinationsrate: 8\% $\to$ 1.2\% (Faithfulness 0.61 $\to$ 0.92, Auto-Retry bei $<$ 0.7)
- Cost/Query: €0.018 $\to$ €0.012 (-34\%, Token-Budget-Check in CI, n=80 Baseline vs. n=80 nach Optimierung)
- Regressions-Erkennung: 2 Min (CI) statt 2 Wochen (User-Beschwerden)
- Model Drift erkannt: gpt-4o-mini Still-Update (2024-07-18) $\to$ Faithfulness -8\% über 14 Tage $\to$ Pinning gpt-4o-mini-2024-07-18 + Canary 5\%
Ohne systematische Evaluation ist jeder Prompt-Change Raten. Du weisst nicht, ob das System besser oder schlechter geworden ist. Du weisst nicht, welche Kosten die Änderung verursacht. Und du weisst nicht, ob die Qualität in Produktion nachlässt.
}

% ============================================================
% 1. WHY THIS MATTERS
% ============================================================
\section{Warum Evaluation die einzige Engineering-Disziplin ist, die probabilistische Systeme beherrschbar macht}

In der klassischen Softwareentwicklung ist Testing Standard: Unit-Tests, Integrationstests, E2E-Tests. Jeder Commit durchläuft eine Pipeline, die Regressionen erkennt. LLMs sind probabilistische Systeme — sie liefern bei identischem Input nicht garantiert denselben Output. Ein Prompt-Update, das einen Fall verbessert, kann zehn andere verschlechtern.

Ohne systematische Evaluation ist jede Prompt-Änderung Raten. Du weisst nicht, ob das System besser oder schlechter geworden ist. Du weisst nicht, welche Kosten die Änderung verursacht. Und du weisst nicht, ob die Qualität in Produktion nachlässt.

Chip Huyen formuliert es präzise: \textbf{Evaluation first} — nicht nachgeschoben, sondern \textbf{vor} RAG, Agents, Fine-Tuning. Evaluation ist die einzige Engineering-Disziplin, die probabilistische Systeme deterministisch beherrschbar macht.

\par
\vspace{0.5cm}
\begin{realitycheckEnv}
Wer "Prompt Engineering = gute Fragen stellen" denkt, baut Systeme, die in Produktion brechen. Prompt Engineering ist: \textbf{Struktur erzwingen, Validierung bauen, Distributionen einengen}. Das ist Engineering. Der Rest ist Hoffen.
\end{realitycheckEnv}

LLMOps — die Übertragung von DevOps-Prinzipien auf LLM-Workflows — schliesst diese Lücke: reproduzierbare Tests, kontinuierliches Monitoring und datengetriebene Entscheidungen über Prompt-Änderungen.

% ============================================================
% 2. MENTAL MODEL
% ============================================================
\section{Mental Model — Drei Dimensionen, ein Trade-off-Dreieck}

Jede LLM-Evaluation misst drei Dimensionen, die in einem harten Trade-off stehen:

\begin{itemize}
    \item \textbf{Qualität} Wie gut ist die Antwort? Richtig, vollständig, sicher, höflich? Gemessen durch Golden Datasets, LLM-Judges oder menschliches Feedback.
    \item \textbf{Kosten} Was kostet ein einzelner Request? Token-Verbrauch (Input + Output), Modell-Wahl, Cache-Trefferrate. Gemessen in Cent pro Query.
    \item \textbf{Latenz} Wie schnell antwortet das System? Time-to-First-Token (TTFT), Total Response Time, P95/P99. Gemessen in Millisekunden.
\end{itemize}

Diese drei Dimensionen stehen im Trade-off: Ein grösseres Modell liefert oft bessere Qualität, kostet aber mehr und ist langsamer. Evaluation hilft dir, den optimalen Punkt für deinen Anwendungsfall zu finden.

\par
\vspace{0.5cm}
\begin{merkeEnv}
\begin{itemize}
    \item LLM $\neq$ deterministische Funktion $\to$ Testing $\neq$ klassisches Unit-Testing
    \item \textbf{Offline} (Golden Set vor Deploy) + \textbf{Online} (Tracing + Drift + A/B in Prod) — keine ersetzt die andere
    \item \textbf{Proxy vs. Business Metrik}: Faithfulness korreliert nicht automatisch mit Customer Satisfaction
    \item Evaluation ist \textbf{kontinuierlicher Prozess}, kein einmaliges Artefakt
\end{itemize}
\end{merkeEnv}

% ============================================================
% 3. ARCHITECTURE
% ============================================================
\section{Architektur einer Evaluations-Pipeline}

Eine vollständige Evaluations-Pipeline besteht aus zwei gekoppelten Ebenen:

\begin{verbatim}
+-------------------------------------------------------------+
|                    CI/CD Pipeline                           |
|  Lint $\to$ Golden Dataset $\to$ Multi-Judge $\to$ Cost/Latency Gate   |
+-------------------------+-----------------------------------+
                          |
+-------------------------v-----------------------------------+
|                  Production Monitoring                       |
|  Tracing (Langfuse) $\to$ Metrics (P95, Cost, Faithfulness)    |
|       $\to$ Alerting (Drift, Cost Spike, Quality Drop)          |
|       $\to$ Dashboard (Grafana/Langfuse) $\to$ A/B Test Framework   |
+-------------------------------------------------------------+
\end{verbatim}

Die Trennung in Offline-Evaluation (vor dem Deployment) und Online-Monitoring (in der Produktion) ist zentral. Keine der beiden Ebenen ersetzt die andere.

\subsection{Komponenten einer Evaluations-Pipeline}

\begin{enumerate}
    \item \textbf{Golden Dataset}: Versionierte Sammlung von Testfällen mit erwarteten Ergebnissen (Git-getaggt, wachsend)
    \item \textbf{Evaluation Runner}: Führt Tests parallel aus, cached Ergebnisse, budget-gated (Max \$5/Run)
    \item \textbf{LLM Judge}: Multi-Judge (3 Modelle, verschiedene Familien, Majority Vote + Agreement-Score)
    \item \textbf{Tracing Layer}: Protokolliert jeden LLM-Aufruf mit Prompt-Hash, Tokens, Latency, Cost, Validation-Status
    \item \textbf{Dashboard + Alerting}: Quality < 0.7, Cost/Query > €0.02, P95 > 2s, Drift (KS/PSI) p < 0.01
\end{enumerate}

\praxishinweis{Ohne Golden Dataset ist jede Prompt-Änderung geraten. Ein Dataset mit 50–100 Testfällen reicht für den Start — aber es muss wachsen. Sammle jeden Produktions-Fehler als neuen Testfall. Die Kombination aus Offline-Evaluation (vor Deployment) und Online-Monitoring (in Produktion) ist Pflicht — keine der beiden Ebenen ersetzt die andere.}

% ============================================================
% 4. CORE CONCEPTS
% ============================================================
\section{Core Concepts — Golden Dataset, LLM Judge, Metriken, Drift, Human Protocol}

\subsection{Golden Dataset — Das Fundament}

Ein Golden Dataset ist eine versionierte Sammlung von Testfällen mit erwarteten Ergebnissen. Struktur:

\begin{itemize}
    \item \textbf{Input}: User-Query (String)
    \item \textbf{Expected Output}: Entweder Keywords (contains), JSON Schema (structured), oder Reference Answer (semantic similarity)
    \item \textbf{Should Not Contain}: Negative Constraints (z.B. "keine Halluzination", "keine PII")
    \item \textbf{Metadata}: Query-Type (Status, Rechnung, Kündigung, Technisch, OOD), Difficulty, Expected Cost Budget
\end{itemize}

\paragraph{Operationalisierung:}
\begin{itemize}
    \item Start: 50–100 Cases, stratifiziert nach Query-Type
    \item Wachstum: Jeder Produktions-Fehler $\to$ neuer Testfall + CI Fail bis Fix
    \item Versionierung: \texttt{golden\_v1.json} $\to$ \texttt{golden\_v2.json} + Git Tag \texttt{eval/golden-v2}
    \item Synthetische Generierung nur als \textbf{Seed}, 20\% manuell geprüft
    \item Train/Test Split strikt getrennt (80/20), Versionierung pro Split
\end{itemize}

\begin{lstlisting}[language=Python, caption={Golden Dataset Schema — [Production Ready]}]
from typing import TypedDict, Literal, Optional
from pydantic import BaseModel, Field

class ExpectedOutput(TypedDict):
    type: Literal["contains", "json_schema", "semantic_similarity"]
    keywords: Optional[list[str]]
    schema: Optional[dict]
    reference_answer: Optional[str]
    min_similarity: Optional[float]

class GoldenCase(BaseModel):
    id: str
    input: str
    query_type: Literal["status", "billing", "cancellation", "technical", "ood"]
    difficulty: Literal["easy", "medium", "hard"]
    expected: ExpectedOutput
    should_not_contain: list[str] = Field(default_factory=list)
    cost_budget_eur: float = 0.02
    latency_budget_ms: int = 2000

# Example
GoldenCase(
    id="status_001",
    input="Wo bleibt meine Bestellung #12345?",
    query_type="status",
    difficulty="easy",
    expected={
        "type": "json_schema",
        "schema": {
            "order_id": str,
            "status": Literal["shipped", "processing", "delivered", "cancelled"],
            "tracking_url": Optional[str],
            "estimated_delivery": Optional[str]
        }
    },
    should_not_contain=["weiss nicht", "keine information", "probier später"],
    cost_budget_eur=0.015,
    latency_budget_ms=1500
)
\end{lstlisting}

\subsection{LLM-as-a-Judge — Bias Taxonomy \& Mitigation}

LLM-as-a-Judge hat einen fundamentalen Bias: \textbf{Das evaluierende Modell bevorzugt oft eigene Output-Charakteristiken}. Ein GPT-4o-Judge bevorzugt lange, ausführliche Antworten. Ein Claude-Judge bevorzugt vorsichtige, abwägende Antworten. Ein Gemini-Judge bewertet strukturierte Aufzählungen höher.

\begin{table}[H]
\centering
\begin{tabularx}{\linewidth}{lXX}
\toprule
\textbf{Bias} & \textbf{Manifestation} & \textbf{Mitigation} \\
\midrule
Verbosity & GPT-4o bevorzugt lange Antworten & Rubric: "Kürze = Bonus", Length-Normalisierung \\
Self-Preference & Modell bewertet eigenen Stil höher & \textbf{Multi-Judge}: 3 verschiedene Modell-Familien \\
Position Bias & Erste Antwort in Pairwise gewinnt & \textbf{Position Swapping}: beide Reihenfolgen testen \\
Style Bias & Claude = vorsichtig, Gemini = strukturiert & Rubric-basiert statt freier Skala \\
Calibration Drift & Judge wird im Laufe der Zeit "nachsichtiger" & Monatliche Human-Calibration (Krippendorff $\alpha \geq 0.8$) \\
\bottomrule
\end{tabularx}
\caption{LLM Judge Bias Taxonomy und Mitigation}
\end{table}

\paragraph{Multi-Judge Implementation:} 3 Modelle (verschiedene Familien: OpenAI, Anthropic, Google) $\to$ Majority Vote + Agreement-Score (Konsens $\geq 66\%$ = valide).

\begin{lstlisting}[language=Python, caption={Multi-Judge mit Majority Vote + Agreement — [Production Ready]}]
from collections import Counter
from dataclasses import dataclass
from typing import Optional

JUDGE_MODELS = [
    "gpt-4o-2024-08-06",      # OpenAI family
    "claude-3-5-sonnet-20241022",  # Anthropic family  
    "gemini-1.5-pro-002"      # Google family
]

@dataclass
class JudgeResult:
    score: Optional[float]
    confidence: float
    agreement: float
    individual_scores: list[int]

def multi_judge(query: str, answer: str, rubric: str) -> JudgeResult:
    """Drei unabhängige Modelle bewerten, Majority-Vote als Ergebnis."""
    scores = []

    for model in JUDGE_MODELS:
        judge_prompt = (
            f"Bewerte die Antwort auf Skala 1-10.\n"
            f"Rubrik: {rubric}\n"
            f"Frage: {query}\n"
            f"Antwort: {answer}\n"
            f"Gib NUR eine Zahl zurück."
        )
        response = call_llm(model, judge_prompt, temperature=0.0)
        try:
            scores.append(int(response.strip()))
        except ValueError:
            continue

    if not scores:
        return JudgeResult(score=None, confidence=0.0, agreement=0.0, individual_scores=[])

    score_counts = Counter(scores)
    majority_score = score_counts.most_common(1)[0][0]
    agreement = max(score_counts.values()) / len(scores)

    return JudgeResult(
        score=majority_score,
        confidence=agreement,
        agreement=agreement,
        individual_scores=scores
    )

# Usage
result = multi_judge(
    query="Was ist der Unterschied zwischen RAG und Fine-Tuning?",
    answer="RAG ergänzt den Prompt mit relevanten Dokumenten, Fine-Tuning passt die Modellgewichte an.",
    rubric="Faktenrichtigkeit, Vollständigkeit, Klarheit"
)
print(f"Score: {result.score}/10 | Agreement: {result.agreement:.0%}")
print(f"Einzelbewertungen: {result.individual_scores}")
\end{lstlisting}

\praxishinweis{Multi-Judge ist Pflicht im CI-Gate. Single-Judge nur f\"ur Exploration. Agreement < 66\% = "uneinheitlich" $\to$ Human Review triggern. Monatliche Kalibrierung gegen Human Ground Truth (Krippendorff $\alpha \geq 0.8$) verhindert Calibration Drift.}

\subsection{RAG-spezifische Metriken (Forward-Ref: Kap. 7)}

Für RAG-Systeme (Kap. 7) brauchst du Metriken, die Retrieval und Generation getrennt messen:

\begin{table}[H]
\centering
\begin{tabularx}{\linewidth}{lX}
\toprule
\textbf{Metrik} & \textbf{Was sie misst} \\
\midrule
Hit Rate @ K & Ist das \textbf{relevante} Chunk in Top-K? (Binary: ja/nein pro Query) \\
MRR @ K & Wie weit oben ist das \textbf{erste} relevante Chunk? \\
Faithfulness & Wie viel der LLM-Antwort ist durch retrievede Chunks \textbf{belegbar}? (RAGAS-Style) \\
Answer Relevance & Beantwortet die Antwort die \textbf{eigentliche} Query? \\
Context Precision & Signal/Rauschen im übergebenen Kontext \\
\bottomrule
\end{tabularx}
\caption{RAG-Metriken — das Minimum für Produktion}
\end{table}

\paragraph{Faithfulness Implementation (RAGAS-Style):}

\begin{lstlisting}[language=Python, caption={Faithfulness Metric — [Production Ready]}]
from sentence_transformers import CrossEncoder
import numpy as np

class FaithfulnessEvaluator:
    """Misst: Wie viel der Answer-Tokens wird durch Retrieved Chunks belegt?"""
    
    def __init__(self, judge_model: str = "gpt-4o-2024-08-06"):
        self.judge_model = judge_model
        self.nli_model = CrossEncoder("cross-encoder/nli-deberta-v3-base")
    
    def score(self, question: str, answer: str, contexts: list[str]) -> float:
        """0.0 = keine Belegung, 1.0 = vollständig belegt."""
        # 1. Answer in Claims zerlegen (einfache Heuristik: Sätze)
        claims = [s.strip() for s in answer.split(".") if s.strip()]
        if not claims:
            return 1.0
        
        # 2. Jeden Claim gegen Contexts prüfen (NLI: Entailment/Contradiction/Neutral)
        supported = 0
        for claim in claims:
            max_entailment = 0.0
            for ctx in contexts:
                pairs = [(claim, ctx)]
                scores = self.nli_model.predict(pairs)
                entailment_prob = scores[0][0]  # Entailment class prob
                max_entailment = max(max_entailment, entailment_prob)
            
            if max_entailment > 0.5:  # Threshold tuneable
                supported += 1
        
        return supported / len(claims)

# Usage
evaluator = FaithfulnessEvaluator()
faithfulness = evaluator.score(
    question="Was ist die Rückgabepolitik?",
    answer="Sie haben 14 Tage Rückgaberecht, kostenloser Versand.",
    contexts=["Rückgaberecht: 14 Tage ab Erhalt. Kostenloser Rückversand."]
)
print(f"Faithfulness: {faithfulness:.2f}")  # Erwartet: ~1.0
\end{lstlisting}

\subsection{Drift Detection — Algorithmen}

Drift ist der stille Killer in Produktion. Drei komplementäre Detektoren:

\paragraph{1. KS-Test (Embedding Distribution Shift):}
Vergleicht die Embedding-Verteilung des aktuellen 7-Tage-Fensters gegen das Referenz-Fenster (Golden-Set-Embeddings). p < 0.01 $\to$ Alert.

\begin{lstlisting}[language=Python, caption={KS-Test auf Embedding Drift — [Production Ready]}]
from scipy import stats
import numpy as np

def detect_embedding_drift(current_embeddings: np.ndarray, 
                           reference_embeddings: np.ndarray,
                           alpha: float = 0.01) -> tuple[bool, float]:
    """
    KS-Test auf jeder Dimension + Bonferroni-Korrektur.
    Returns: (drift_detected, min_p_value)
    """
    n_dims = current_embeddings.shape[1]
    p_values = []
    
    for dim in range(n_dims):
        stat, p = stats.ks_2samp(
            current_embeddings[:, dim], 
            reference_embeddings[:, dim]
        )
        p_values.append(p)
    
    # Bonferroni: alpha / n_dims
    corrected_alpha = alpha / n_dims
    drift = any(p < corrected_alpha for p in p_values)
    return drift, min(p_values)

# Production: Wöchentlich auf letzten 7 Tage vs. Golden Set Reference
# Alert bei drift=True $\to$ Model Pinning Review triggern
\end{lstlisting}

\paragraph{2. PSI (Population Stability Index) — Kosten/Latency/Token-Länge:}
Misst Verteilungs-Shift bei skalaren Metriken. PSI > 0.2 = signifikanter Shift.

\begin{lstlisting}[language=Python, caption={PSI für Kosten/Latency Drift — [Production Ready]}]
def psi(expected: np.ndarray, actual: np.ndarray, buckets: int = 10) -> float:
    """Population Stability Index. > 0.2 = signifikanter Shift."""
    # Quantile-basierte Buckets auf Expected
    breakpoints = np.percentile(expected, np.linspace(0, 100, buckets + 1))
    breakpoints[0] = -np.inf
    breakpoints[-1] = np.inf
    
    exp_perc = np.histogram(expected, bins=breakpoints)[0] / len(expected)
    act_perc = np.histogram(actual, bins=breakpoints)[0] / len(actual)
    
    # Vermeide Division durch Null
    exp_perc = np.where(exp_perc == 0, 0.0001, exp_perc)
    act_perc = np.where(act_perc == 0, 0.0001, act_perc)
    
    return np.sum((act_perc - exp_perc) * np.log(act_perc / exp_perc))

# Monitoring: Wöchentlich PSI(Cost/Query) vs. Baseline (Golden Set Runs)
# Alert bei PSI > 0.2
\end{lstlisting}

\paragraph{3. Quality Drift (Faithfulness/Relevance Trend):}
Lineare Regression über 7-Tage-Fenster. Steigung < -5\% $\to$ Alert.

\begin{lstlisting}[language=Python, caption={Quality Drift Detection — [Production Ready]}]
from scipy import stats

def detect_quality_drift(daily_scores: list[float], 
                         window_days: int = 7,
                         threshold_slope: float = -0.05) -> tuple[bool, float]:
    """
    daily_scores: Liste täglicher Faithfulness/Relevance Scores (neueste zuletzt)
    Returns: (drift_detected, slope_per_day)
    """
    if len(daily_scores) < window_days:
        return False, 0.0
    
    recent = daily_scores[-window_days:]
    x = np.arange(len(recent))
    slope, intercept, r_value, p_value, std_err = stats.linregress(x, recent)
    
    drift = slope < threshold_slope and p_value < 0.05
    return drift, slope

# Alert: drift=True + slope < -0.05/Tag $\to$ Model Pinning Review + Canary
\end{lstlisting}

\subsection{Human Evaluation Protocol}

LLM-Judges ersetzen keine domänenspezifische menschliche Bewertung. Protokoll:

\begin{itemize}
    \item \textbf{Annotation Guidelines}: 3-seitiges PDF mit Good/Bad/Edge-Beispielen
    \item \textbf{Krippendorff's $\alpha$}: Inter-Annotator-Agreement Ziel $\geq 0.8$
    \item \textbf{Sample Size}: Min 50 Cases/Woche, stratifiziert nach Query-Type
    \item \textbf{Blind}: Annotatoren sehen nicht, welcher Prompt/Model dahintersteckt
    \item \textbf{Frequenz}: Wöchentlich 50 Cases (nicht monatlich 200)
\end{itemize}

\begin{lstlisting}[language=Python, caption={Krippendorff Alpha Berechnung — [Didaktisch vereinfacht]}]
import krippendorff

# data: [annotator1_scores, annotator2_scores, annotator3_scores]
# Jede Liste: Scores für dieselben 50 Cases (1-10)
alpha = krippendorff.alpha(reliability_data=data, level_of_measurement="interval")
print(f"Krippendorff \alpha: {alpha:.3f}")  # Ziel: \geq 0.80

# \alpha < 0.67 = unzuverlässig, 0.67-0.80 = tentativ, \geq 0.80 = verlässlich
# Bei niedrigem \alpha: Guidelines schärfen, Annotatoren schulen, Cases klären
\end{lstlisting}

% ============================================================
% 5. PRODUCTION EXAMPLE — SUPPORTPILOT EVALUATION JOURNEY
% ============================================================
\section{Production Example — SupportPilot Evaluation Journey}

\textbf{Ausgangslage}: Prompt-v1 deployed, "fühlt sich gut an", aber keine Metriken.

\paragraph{Schritt 1 — Golden Set Aufbau (Woche 1):}
\begin{itemize}
    \item 80 reale User-Queries aus Logs gesampelt (stratifiziert: 30\% Status, 25\% Rechnung, 20\% Kündigung, 15\% Technisch, 10\% OOD)
    \item Manuelle Annotation: Expected Keywords + JSON Schema + "Should Not Contain"
    \item Baseline: Prompt-v1 $\to$ Pass Rate 72\%, Avg Faithfulness 0.61, Cost/Query €0.018, P95 1.8s
\end{itemize}

\paragraph{Schritt 2 — CI-Gate etabliert (Woche 2):}
\begin{itemize}
    \item GitHub Action: \texttt{pytest eval/golden.py --threshold 0.85}
    \item Multi-Judge (GPT-4o, Claude 3.5 Sonnet, Gemini 1.5 Pro) $\to$ Majority Vote
    \item Budget Gate: Max \$5/Run, Cache 80\% Savings
\end{itemize}

\paragraph{Schritt 3 — Regression gefangen (Woche 3):}
\begin{itemize}
    \item Prompt-v2 (neue Few-Shot Examples) $\to$ CI: Pass Rate 78\% (grün), aber Faithfulness 0.58 (rot)
    \item Auto-Rollback verhindert Deploy
    \item Root Cause: Neues Example halluzinierte "30 Tage Rückgabe" statt "14 Tage"
\end{itemize}

\paragraph{Schritt 4 — Production Drift erkannt (Monat 2):}
\begin{itemize}
    \item Langfuse Alert: Faithfulness Trend -8\% über 14 Tage
    \item Ursache: Provider hob gpt-4o-mini stillschweigend auf neue Version (2024-07-18)
    \item Fix: Model Pinning \texttt{gpt-4o-mini-2024-07-18} + Canary 5\% Traffic
\end{itemize}

\paragraph{Ergebnis nach 3 Monaten:}
\begin{itemize}
    \item Pass Rate: 72\% $\to$ 94\% (Golden Set, n=80, Agreement $\geq 66\%$)
    \item Halluzinationsrate: 8\% $\to$ 1.2\% (Faithfulness 0.61 $\to$ 0.92, Auto-Retry bei < 0.7)
    \item Cost/Query: €0.018 $\to$ €0.012 (-34\%, Token-Budget-Check in CI, n=80 Baseline vs. n=80 nach Optimierung)
    \item Regressions-Erkennung: 2 Min (CI) statt 2 Wochen (User-Beschwerden)
    \item Model Drift erkannt: gpt-4o-mini Still-Update (2024-07-18) $\to$ Faithfulness -8\% über 14 Tage $\to$ Pinning + Canary 5\%
\end{itemize}

\paragraph{Messrahmen für obige Zahlen:}
\begin{itemize}
    \item \textbf{Baseline}: Prompt-v1, n=80 Golden Cases, 3 Judge-Modelle, Agreement $\geq 66\%$
    \item \textbf{Evaluation Method}: Offline Golden Set + Online 10\% Sample Human Review
    \item \textbf{Dataset}: SupportPilot Production Queries Jan–Mar 2024, stratified
    \item \textbf{n}: 80 (Golden) + 200/Woche (Online Sample)
    \item \textbf{Metric}: Pass Rate (Binary), Faithfulness (0-1), Cost/Query (EUR), P95 Latency (ms)
    \item \textbf{What Changed}: Golden Set + CI Gate + Multi-Judge + Model Pinning + Drift Alerts
    \item \textbf{Limitations}: Kein A/B Test (Traffic zu niedrig), Human Review nur 10\%, Judge-Kalibrierung monatlich
\end{itemize}

% ============================================================
% 6. TRADE-OFFS
% ============================================================
\section{Trade-offs}

\begin{table}[H]
\centering
\begin{tabularx}{\linewidth}{lXXX}
\toprule
\textbf{Entscheidung} & \textbf{Option A} & \textbf{Option B} & \textbf{Empfehlung} \\
\midrule
Judge-Count & Single (billig) & Multi-Judge 3x (teuer) & \textbf{Multi-Judge} für CI-Gate; Single für Exploration \\
Metric Type & Proxy (Faithfulness) & Business (CSAT) & \textbf{Beide}: Proxy im CI, Business im Monthly Review \\
Dataset Size & 50 (schnell) & 500+ (robust) & \textbf{Start 50}, wöchentlich +5 aus Prod-Fehlern \\
Human Eval & Wöchentlich 50 & Monatlich 200 & \textbf{Wöchentlich 50} stratifiziert; $\alpha \geq 0.8$ prüfen \\
Drift Sensitivity & KS-Test (sensitiv) & PSI (robust) & \textbf{KS für Embeddings}, PSI für Kosten/Latency \\
A/B Test & Full Randomisierung & Canary (5\%) & \textbf{Canary für Prompt-Changes}, Full A/B für Model Swaps \\
\bottomrule
\end{tabularx}
\caption{Trade-offs in der Evaluation — Entscheidungen mit Begründung}
\end{table}

% ============================================================
% 7. FAILURE MODES
% ============================================================
\section{Failure Modes — Was schiefgeht und wie du es verhinderst}

\begin{table}[H]
\centering
\begin{tabularx}{\linewidth}{lXXX}
\toprule
\textbf{Failure Mode} & \textbf{Symptom} & \textbf{Detection} & \textbf{Prevention} \\
\midrule
Judge Bias nicht kalibriert & Judge-Score $\uparrow$, Human-Score $\downarrow$ & Monthly Krippendorff $\alpha < 0.8$ & Monatliche Human-Calibration \\
Golden Set veraltet & Pass Rate $\uparrow$, Prod-Quality $\downarrow$ & Online/Offline Gap > 10\% & Jeder Prod-Fehler $\to$ Golden Set + CI Fail \\
Dataset Leakage & Perfekt auf Golden, Müll in Prod & Train/Test Split verwechselt & Streng getrennte Splits, Versionierung \\
Metric Gaming & Optimierung auf Faithfulness $\to$ Answer Relevance $\downarrow$ & Guardrail-Metriken im Dashboard & \textbf{Immer} Primär + Sekundär + Guardrail \\
Cost Explosion in Eval & \$50/Run statt \$5 & Budget Gate in CI & Max \$5/Run, Cache 80\%, Delta-Eval \\
Alert Fatigue & Drift-Alerts täglich, keiner schaut & Alert/Week > 3 ohne Action & SLO-basiert: nur Alert bei SLO-Verletzung \\
\bottomrule
\end{tabularx}
\caption{Top 6 Failure Modes in der Evaluation}
\end{table}

% ============================================================
% 8. EVALUATION — META-EVALUATION
% ============================================================
\section{Meta-Evaluation — Wie evaluiere ich meine Evaluation?}

\begin{itemize}
    \item \textbf{Offline/Online Correlation}: Pearson r zwischen Golden-Set-Pass-Rate und Prod-Quality (Ziel: r > 0.7)
    \item \textbf{Judge Calibration}: Monthly Human vs. Judge Agreement (Krippendorff $\alpha$)
    \item \textbf{False Positive Rate}: Alerts fired / True Regressions (Ziel: < 30\%)
    \item \textbf{Time to Detect}: Median Zeit von Regressions-Einführung bis Alert (Ziel: < 10 Min)
    \item \textbf{Cost per Detection}: Eval-Kosten / gefundene Regression (Ziel: < \$10)
\end{itemize}

\praxishinweis{Messe deine Evaluation. Wenn du nicht weisst, wie gut dein Eval-System funktioniert, fliegst du blind — genau das, was du vermeiden willst.}

% ============================================================
% 9. BEST PRACTICES
% ============================================================
\section{Best Practices für LLM-Evaluation}

\begin{enumerate}
    \item \textbf{Golden Dataset versionieren} — Git Tags \texttt{eval/golden-v\{N\}}, Changelog pro Version
    \item \textbf{CI-Gate hart} — Exit Code 1 bei Pass Rate < 85\% ODER Faithfulness < 0.8 ODER Cost > Budget
    \item \textbf{Multi-Judge Pflicht} — 3 Modelle, verschiedene Familien, Majority Vote + Agreement
    \item \textbf{Budget Gate in CI} — Max \$5/Run, Delta-Eval für geänderte Prompts
    \item \textbf{Prompt-Hash in Traces} — Jeder Trace verknüpft Prompt-Version $\to$ Golden-Set-Run
    \item \textbf{Synthetisch $\to$ Human Pipeline} — LLM generiert Test-Cases $\to$ Human validiert $\to$ Golden Set
    \item \textbf{Drift Detection Dual} — KS-Test (Embeddings) + PSI (Cost/Latency) + Quality Trend
    \item \textbf{Human Evaluation Protocol} — Guidelines, Krippendorff $\alpha \geq 0.8$, Blind, Stratified
    \item \textbf{A/B Test Framework} — Canary 5\%, Minimum Detectable Effect definiert, Power Analysis
    \item \textbf{Documentation} — Jeder Metrik: Definition, Baseline, Threshold, Owner, Last Reviewed
\end{enumerate}

% ============================================================
% 10. ANTI-PATTERNS
% ============================================================
\section{Anti-Patterns — Was du vermeiden solltest}

\begin{enumerate}
    \item \textbf{Evaluation ohne Golden Dataset} $\to$ LLM-Judge misst nur Judge-Meinung
    \item \textbf{Judge = zu evaluierendes Modell} $\to$ Selbstbestätigung, keine externe Validierung
    \item \textbf{Nur Accuracy messen} $\to$ System 100\% Accuracy, aber 10s Latency / \$1/Query / unsicher
    \item \textbf{Einmaliges Dataset} $\to$ Nach 3 Monaten irrelevant (User-Verhalten, Model-Updates, Produkt)
    \item \textbf{Human Evaluation ignorieren} $\to$ LLM-Judge ersetzt keine Domänen-Expertise (Medizin, Recht, Finance)
    \item \textbf{Online Monitoring vergessen} $\to$ Offline grün, Prod rot (Model Drift, Data Drift, Provider Drift)
    \item \textbf{Single Metric Optimization} $\to$ Faithfulness $\uparrow$ aber Answer Relevance $\downarrow$ = nutzlose Antworten
    \item \textbf{Kein Budget Gate} $\to$ \$500 Eval-Run blockiert CI, Team deaktiviert Tests
\end{enumerate}

% ============================================================
% 11. PRODUCTION CHECKLIST
% ============================================================
\section{Production Checklist}

\begin{itemize}
    \item [ ] Golden Dataset $\geq 50$ Cases, versioniert in Git, wächst wöchentlich
    \item [ ] CI Pipeline: Lint $\to$ Golden Set $\to$ Multi-Judge $\to$ Cost/Latency Gate $\to$ Exit Code
    \item [ ] Multi-Judge: 3 Modelle (verschiedene Familien), Agreement $\geq 66\%$
    \item [ ] Tracing: Jeder Call geloggt (Prompt-Hash, Tokens, Latency, Cost, Validation)
    \item [ ] Dashboard: Quality (Faithfulness/Relevance), Cost/Query, P95 Latency, Error Rate
    \item [ ] Alerts: Quality < 0.7, Cost > 2x Baseline, P95 > 2s, Drift (KS/PSI) p < 0.01
    \item [ ] Model Pinning: Explizite Versionen (\texttt{gpt-4o-mini-2024-07-18}), Canary 5\%
    \item [ ] Human Eval: Wöchentlich 50 Cases, Krippendorff $\alpha \geq 0.8$, Guidelines dokumentiert
    \item [ ] A/B Framework: Canary für Prompts, Full Randomisierung für Model Swaps
    \item [ ] Runbooks: Quality Drop, Cost Spike, Provider Outage, Drift Alert
\end{itemize}

% ============================================================
% 12. EXERCISES
% ============================================================
\section{Übungen}

\begin{enumerate}
    \item \textbf{Golden Set Bootstrap}: Erstelle 30 Cases für SupportPilot (10 Status, 8 Rechnung, 7 Kündigung, 5 Technisch) mit Expected Keywords + JSON Schema
    \item \textbf{Multi-Judge Implementation}: Implementiere Majority Vote + Agreement für 3 Judge-Modelle
    \item \textbf{CI Gate}: GitHub Action mit Pass Rate $\geq 85\%$, Faithfulness $\geq 0.8$, Cost $\leq \$5$
    \item \textbf{Drift Alert}: KS-Test auf Embedding-Distribution + PSI auf Cost/Query, Alert bei p < 0.01
    \item \textbf{Human Evaluation}: Annotation Guidelines für 3 Annotatoren, Krippendorff $\alpha$ Berechnung
    \item \textbf{A/B Test Design}: Sample Size Calculator für Prompt-v1 vs v2 (MDE 5\%, Power 80\%, $\alpha$ 5\%)
    \item \textbf{Cost Anomaly Detection}: Rolling 7-Day Cost/Query, Alert bei +50\% vs. Baseline
\end{enumerate}

% ============================================================
% 13. FURTHER READING
% ============================================================
\section{Weiterführende Ressourcen}

\begin{itemize}
    \item \textbf{Primary}: Chip Huyen — "AI Engineering", Kap. 3 (Evaluation); "Designing ML Systems", Kap. 8
    \item \textbf{Papers}: 
        \begin{itemize}
            \item Liu et al. (2023) — "LLM-as-a-Judge" (Bias Taxonomy)
            \item Sarti et al. (2024) — "RAGAS: Automated Evaluation of RAG"
            \item Yan et al. (2024) — "Corrective RAG" (CRAG Quality Gate)
            \item Zheng et al. (2024) — "Judging LLM-as-a-Judge with MT-Bench and Chatbot Arena"
        \end{itemize}
    \item \textbf{Tools}: Langfuse (Open Source), DeepEval, Promptfoo, RAGAS, OpenAI Evals, Garak (Security)
    \item \textbf{Standards}: ISO/IEC 25059 (AI Quality Model), NIST AI RMF (Evaluation)
\end{itemize}

% ============================================================
% TEASER FÜR NÄCHSTES KAPITEL
% ============================================================

\par\mbox{}\par
\begin{AUTORnotizEnv}
Nächstes Kapitel: RAG und Vector Datenbanken. Da lernst du: Chunking als Hebel (Recursive 512/64), Hybrid Search (BM25 + Embedding), Cross-Encoder Re-Ranker, Faithfulness-Metrik für Generation. Die Evaluation-Disziplin hier zahlt sich aus: Golden Set für Retrieval (Hit Rate @ 5, MRR), Faithfulness für Generation. Ohne Eval-Gate fliegst du blind in RAG.
\end{AUTORnotizEnv}

\textbf{Nächster Schritt:} Jetzt weisst du, wie du misst. Im nächsten Kapitel geht es um RAG und Vector Datenbanken — wie du LLMs mit deinem eigenen Wissen anreicherst. Die Evaluation-Disziplin hier zahlt sich aus: Golden Set für Retrieval, Faithfulness-Metrik für Generation.
\endinput

\clearpage